\documentclass[conference,onecolumn]{IEEEtran}

\usepackage[T1]{fontenc}
\usepackage[utf8]{inputenc}
\usepackage[spanish,es-tabla]{babel}
\usepackage{lmodern}
\usepackage{microtype}
\usepackage{booktabs}
\usepackage{float}
\usepackage{array}
\usepackage{longtable}
\usepackage{url}
\usepackage[hidelinks]{hyperref}

\newcolumntype{L}[1]{>{\raggedright\arraybackslash}p{#1}}

\title{EV-PRAC-01: Evaluación de la Unidad I}
\author{\IEEEauthorblockN{Danela Dayana Arteaga Alava}\IEEEauthorblockA{Carrera de Software\\
Facultad de Ciencias de la Computación\\
Universidad Técnica Estatal de Quevedo, Ecuador\\
Ingeniería de Requerimientos -- ISR-401}}

\begin{document}
\maketitle

\noindent\emph{Esta es la transcripción fiel de la prueba escrita a mano el 18 de noviembre de 2026; las fotografías de las hojas están en la carpeta evidencias/ del repositorio.}

\bigskip

% ===================== Hoja 1 =====================
\noindent\textbf{Hoja 1}

\noindent Nombre: Danela Dayana Arteaga Alava\\
Fecha: 18 de Noviembre del 2026\\
Curso: 4 to Paralelo “B”

\begin{center}
Evaluación de la Unidad I (EV-PRAC-01)
\end{center}

\section*{T1.}

\noindent a) Explique con sus palabras

Según ISO/IEC/IEEE 29148 un requisito es una necesidad o condición que deberá cumplir un sistema para satisfacer lo necesitado o acordado con los interesados.

Decisión de si E7 y E12 son requisitos del producto ALERTARIO

\noindent E7: No es un requisito del producto, ya que solo describe el procedimiento que actualmente se utiliza para medir.

\noindent E12: No es un requisito de producto, es un requisito de proceso, ya que este nos describe como se realizarán las reuniones y quién firmará el acta de requisitos.

\medskip
\noindent b) Explique qué describe la especificación de requisitos de los interesados.

\noindent StRS: Describe las necesidades y expectativas de los interesados, por ejemplo los presidentes barriales, ya que ellos necesitan alertas.

\noindent SyRS: Describe los requisitos del sistema, incluyendo hardware, software e interfaces, por ejemplo las estaciones HL-20 que envian lecturas con el protocolo LoRaWAN, de eso el interesado es el técnico de TIC

\noindent SRS: Describe las funciones y caracteristicas que debe cumplir el software, por ejemplo el jefe de la unidad de GR, ya que genera informes mensuales en formato PDF

\medskip
\noindent c) Para cada hecho H1 a H5

\begin{table}[H]
\centering
\begin{tabular}{@{}lll@{}}
\toprule
\textbf{Hecho} & \textbf{Proceso} & \textbf{Pohl}\\
\midrule
H1 & Elicitación & Central\\
H2 & Analisis y negociación & Central\\
H3 & Especificación & Central\\
H4 & Validación & Central\\
H5 & Gestión & Transversal\\
\bottomrule
\end{tabular}
\end{table}

% ===================== Hoja 2 =====================
\noindent\textbf{Hoja 2}

\section*{T2.}

\noindent a) Clasificación. Dibuje una tabla para E1 a E12

{\small
\begin{longtable}{@{}L{0.06\textwidth}L{0.24\textwidth}L{0.16\textwidth}L{0.46\textwidth}@{}}
\toprule
\textbf{ID} & \textbf{TIPO} & \textbf{DOCUMENTO} & \textbf{JUSTIFICACIÓN}\\
\midrule
\endhead
E1 & Funcional & StRS & Expresa la necesidad de recibir alertas\\
E2 & Requisito de dominio & SyRS & Establece los umbrales de protocolo\\
E3 & Restricción & SyRS & Define la frecuencia de comunicación y el protocolo\\
E4 & Restricción & -- & Solo establece plazo y presupuesto.\\
E5 & Funcional & SRS & Define la generación de documentos PDF.\\
E6 & No funcional & SyRS & Exige continuidad durante los cortes de energia\\
E7 & No es ur requisito & -- & Describe un proceso actual\\
E8 & No funcional & SRS & Establece el tiempo de carga\\
E9 & Restricción tecnologica & SyRS & Determina donde se realizara la instalación y cual sistema operativo usará\\
E10 & Funcional & SRS & Establece permisos y acciones.\\
E11 & No funcional & SRS & Establece tiempo max. de registro.\\
E12 & Requisito de proceso & -- & Describe y define las reuniones y la firma de documentos.\\
\bottomrule
\end{longtable}
}

\noindent b) Auditoria. Dibuje una tabla para B1 a B6

% Las filas B3 a B6 continúan en la hoja 3.
{\small
\begin{longtable}{@{}L{0.05\textwidth}L{0.16\textwidth}L{0.22\textwidth}L{0.13\textwidth}L{0.36\textwidth}@{}}
\toprule
\textbf{ID} & \textbf{TERMINO} & \textbf{CARACTERISTICA} & \textbf{SMART} & \textbf{JUSTIFICACIÓN}\\
\midrule
\endhead
B1 & veloz e intuitivo & No ambiguo y verificable & Medible & No establece valores.\\
B2 & suficiente anticipacion & No ambiguo y verificable. & Medible & No define cuanto tiempo es suficiente\\
B3 & todos los puntos, todo momento & Factible y no ambiguo & Realizable & Exige condiciones imposibles\\
B4 & Ninguno & Esta bien formado & Trazable & Es especifico y verificable, pero no hay justificacion\\
B5 & En lo posible & Singular, no ambiguo y verificable & Especifico & No determina formatos, ni condiciones\\
B6 & Ninguno & Bien formado & Cumple & Define condición, accion y plazos\\
\bottomrule
\end{longtable}
}

% ===================== Hoja 3 =====================
\noindent\textbf{Hoja 3}

\section*{T3.}

\noindent Redacte seis requisitos bien formados

% La fila R6 continúa en la hoja 4.
{\small
\setlength{\tabcolsep}{3pt}
\begin{longtable}{@{}L{0.06\textwidth}L{0.38\textwidth}L{0.11\textwidth}L{0.11\textwidth}L{0.12\textwidth}L{0.14\textwidth}@{}}
\toprule
\textbf{ID} & \textbf{ENUNCIADO} & \textbf{TIPO} & \textbf{DOCUMENTO} & \textbf{FUENTE} & \textbf{M. V.}\\
\midrule
\endhead
B1.a & Cuando un vecino inicie su registro, el sistema deberá permitir completar el registro para recibir alertas en un maximo de 3 minutos & No funcional & SRS & B1 - E11 & Prueba.\\
B1.b & Cuando existan 500 usuarios conectados simultaneamente, el sistema deberá cargar el mapa de las seis estaciones en menos de 3 segundos & No funcional & SRS & B1 - E8 & Prueba\\
B2 & Cuando el nivel del Río supere los 5,5 metros en el puente principal, el sistema deberá enviar un SMS a todos los vecinos registrados en un maximo de 2 segundos & Funcional & SRS & B2 - E2 - B6 & Prueba\\
B3 & Durante el funcionamiento normal y cortes electricos de hasta 12 horas, el sistema deberá registrar las lecturas de las seis estaciones cada 5 minutos & Funcional & SyRS & B3 - E3 - E6 & Prueba\\
B5 & Al finalizar cada mes, el sistema deberá permitir descargar en formato Excel los datos de las 6 estaciones & Funcional & SRS & B5, E5 & Demostración\\
R6 & Cuando un operador autorizado solicite cancelar una alerta, el sistema deberá permitir cancelar y registrar el usuario responsable y la fecha y hora de la acción & Funcional & SRS & E10 - H4 & Prueba\\
\bottomrule
\end{longtable}
}

\end{document}
